\section{Pruebas y Validación}

Si el diseño describe el plan y la implementación lo ejecuta, las pruebas son el momento de la verdad: el instante en que el sistema se somete a condiciones reales y se mide si el comportamiento observado coincide con el esperado. Este capítulo documenta la metodología, los protocolos, los escenarios y los resultados de las pruebas realizadas sobre los cuatro objetivos específicos del proyecto. La documentación exhaustiva de cada protocolo garantiza que los resultados sean reproducibles por cualquier investigador que disponga del mismo testbed, y que las conclusiones sean atribuibles a los CNIs y no al procedimiento de medición.

Un principio fundamental que guió todas las pruebas fue el control de variables: en cada experimento, se modificó únicamente el CNI bajo evaluación, manteniendo constantes el hardware, el sistema operativo, la versión de Kubernetes, el tamaño de los paquetes, la duración de las pruebas y las condiciones de red. Este control es lo que permite atribuir las diferencias en los resultados al CNI y no a factores externos.

\subsection{Metodología General de Pruebas}

\subsubsection{Control de Condiciones: el Escenario de Referencia (Baseline)}

Se definió un escenario de referencia con los siguientes parámetros fijos para todos los CNIs:

\begin{table*}[htbp]
\centering
\scriptsize
\begin{tabularx}{\textwidth}{|Y|Y|}
\hline
\textbf{Parámetro Controlado} & \textbf{Valor Fijo para Todo el Experimento} \\
\hline
Tipo de instancia cloud & DigitalOcean Droplet --- Ubuntu 22.04 LTS, 2 vCPU, 4 GB RAM para nodos Worker \\
Versión de Kubernetes & K3s canal stable --- la misma versión en todos los ciclos de prueba \\
Región del datacenter & fra1 (Frankfurt) --- mismo datacenter para todos los nodos en todas las pruebas \\
Red entre nodos & VPC privada DigitalOcean --- MTU 1500 bytes, sin tráfico externo \\
Duración de cada prueba iperf3 & 300 segundos (5 minutos) por corrida \\
Número de corridas por CNI & 5 corridas independientes --- permite calcular promedios y detectar variabilidad \\
Protocolo de red bajo prueba & TCP --- el protocolo dominante en aplicaciones empresariales reales \\
Herramienta de throughput & iperf3 con salida JSON (\texttt{-J}), un solo stream TCP (\texttt{-P 1}) \\
Herramienta de latencia & Script Bash con \texttt{nc -z -w 2} --- 30 muestras por corrida \\
\hline
\end{tabularx}
\caption{Parámetros del escenario de referencia (baseline) para todas las pruebas.}
\end{table*}

\subsubsection{Automatización: de la Prueba Manual a la Prueba Orquestada}

Todas las pruebas se orquestaron mediante CronJobs de Kubernetes. La secuencia de ejecución automática en cada ciclo fue:
\begin{enumerate}
\item CronJob del cliente iperf3 se activa (minutos 0 y 30). Kubernetes crea el Pod cliente en un nodo diferente al servidor (podAntiAffinity suave, weight 100).
\item Pod cliente ejecuta iperf3 durante 300 segundos y guarda el resultado JSON.
\item CronJob del cliente de latencia ejecuta 30 mediciones de TCP connect (minutos 10 y 40) y guarda su resultado JSON.
\item Al finalizar ambas pruebas, el Resource Usage Exporter (minutos 25 y 55) consulta Prometheus para extraer el consumo de CPU y RAM del agente CNI.
\item \texttt{procesador.js} actualiza \texttt{cni-data.json} con los nuevos datos procesados.
\end{enumerate}

La orquestación automática elimina la variabilidad humana de la ejecución: cada corrida sigue exactamente la misma secuencia con los mismos intervalos, de modo que las diferencias observadas entre CNIs se deben al plugin y no al operador.

\begin{quote}
\textbf{¿Por qué 5 corridas y no 1 o 100?}
Con 1 corrida, un único evento anómalo puede dominar completamente el resultado. Con 100 corridas, el tiempo de prueba se vuelve imprácticamente largo. Con 5 corridas bien distribuidas en el tiempo, el algoritmo IQR puede identificar y descartar la corrida anómala si ocurre, y el promedio de las restantes es estadísticamente representativo. Este balance entre rigor estadístico y viabilidad operativa es estándar en la investigación de redes en entornos cloud \cite{ref7,ref8}.
\end{quote}

\subsection{Pruebas de Rendimiento de Red --- Objetivo Específico 1}

\subsubsection{Protocolo de Prueba de Throughput con iperf3}

iperf3 establece una conexión entre un proceso servidor y un proceso cliente, y el cliente envía datos lo más rápido que la red permite durante el tiempo configurado. Se usó el parámetro \texttt{-P 1} (un solo stream TCP) para representar el comportamiento de una aplicación empresarial típica, no la carga máxima teórica.

\begin{lstlisting}
# Comando iperf3 ejecutado por el CronJob cliente
# -c: servidor  -t: 300s  -J: JSON  -P 1: un solo stream
iperf3 -c iperf-server-svc.cni-benchmark.svc.cluster.local \
  -t 300 -J -P 1 \
  > /results/throughput-$(date +%Y%m%d_%H%M%S)-${CNI_NAME}.json
# El output JSON incluye:
# - bits_per_second (sender): throughput desde el emisor
# - bits_per_second (receiver): throughput real recibido
# - retransmits: paquetes TCP que debieron reenviarse
# - mean_rtt: RTT promedio en microsegundos
\end{lstlisting}

\subsubsection{Resultados de Throughput TCP por CNI}

Los resultados presentados son los promedios de las 5 corridas de 300 segundos para cada CNI, después de aplicar el filtro IQR para eliminar outliers.

\begin{table*}[htbp]
\centering
\scriptsize
\begin{tabularx}{\textwidth}{|Y|Y|Y|Y|Y|}
\hline
\textbf{Métrica} & \textbf{Flannel (VXLAN)} & \textbf{Calico (VXLAN)} & \textbf{Cilium (eBPF)} & \textbf{Antrea (OVS)} \\
\hline
Throughput Sender promedio & 9.21 Gbps & 9.48 Gbps & 9.67 Gbps & 9.35 Gbps \\
Throughput Receiver promedio & 9.18 Gbps & 9.45 Gbps & 9.64 Gbps & 9.31 Gbps \\
Retransmisiones TCP (promedio) & 12.4 / sesión & 4.8 / sesión & 2.1 / sesión & 7.3 / sesión \\
Diferencia Sender-Receiver & $\sim$0.03 Gbps & $\sim$0.03 Gbps & $\sim$0.03 Gbps & $\sim$0.04 Gbps \\
Variabilidad entre corridas ($\sigma$) & $\pm$0.18 Gbps & $\pm$0.09 Gbps & $\pm$0.07 Gbps & $\pm$0.12 Gbps \\
Outliers descartados por IQR & 1 de 5 (20\%) & 0 de 5 (0\%) & 0 de 5 (0\%) & 1 de 5 (20\%) \\
\hline
\end{tabularx}
\caption{Resultados de throughput TCP por CNI (promedio de 5 corridas de 300 segundos).}
\end{table*}

Todos los CNIs se acercan al límite teórico de la red VPC de DigitalOcean ($\sim$9.7--9.8 Gbps). Sin embargo, las diferencias en retransmisiones son significativas: Cilium con eBPF genera apenas 2.1 retransmisiones por sesión, mientras que Flannel genera 12.4, casi 6 veces más. Esta diferencia se explica porque eBPF opera directamente en el kernel y puede tomar decisiones de enrutamiento mucho más rápido que el mecanismo VXLAN de Flannel.

\subsubsection{Protocolo de Prueba de Latencia TCP}

La latencia de establecimiento TCP mide el tiempo que tarda en completarse el handshake de tres vías de TCP, que ocurre cada vez que una aplicación inicia una nueva conexión. El cliente ejecuta \texttt{nc -z -w 2} al servicio iperf3-server, mide el tiempo en ms y calcula min/avg/max/mdev sobre 30 muestras por corrida. Esta métrica es especialmente relevante para arquitecturas de microservicios, donde el costo de iniciar una conexión puede dominar el tiempo total de respuesta y el throughput del enlace no captura esa penalización.

\subsubsection{Resultados de Latencia TCP por CNI}

\begin{table*}[htbp]
\centering
\scriptsize
\begin{tabularx}{\textwidth}{|Y|Y|Y|Y|Y|}
\hline
\textbf{Métrica de Latencia} & \textbf{Flannel (VXLAN)} & \textbf{Calico (VXLAN)} & \textbf{Cilium (eBPF)} & \textbf{Antrea (OVS)} \\
\hline
Latencia mínima (ms) & 0.412 ms & 0.287 ms & 0.198 ms & 0.341 ms \\
Latencia promedio (ms) & 0.687 ms & 0.431 ms & 0.312 ms & 0.512 ms \\
Latencia máxima (ms) & 2.341 ms & 0.891 ms & 0.543 ms & 1.124 ms \\
Jitter / Desviación estándar (ms) & $\pm$0.389 ms & $\pm$0.142 ms & $\pm$0.087 ms & $\pm$0.201 ms \\
Percentil P95 (ms) & 1.834 ms & 0.712 ms & 0.478 ms & 0.893 ms \\
Percentil P99 (ms) & 2.156 ms & 0.843 ms & 0.521 ms & 1.067 ms \\
\hline
\end{tabularx}
\caption{Resultados de latencia TCP por CNI (30 muestras $\times$ 5 corridas, outliers IQR eliminados).}
\end{table*}

Los datos de latencia revelan la mayor brecha de rendimiento entre CNIs. Cilium logra una latencia promedio de 0.312 ms, mientras que Flannel alcanza 0.687 ms, el doble. El P99 de Flannel (2.156 ms) significa que el 1\% de las conexiones tardan más de 2 ms en establecerse. En una aplicación que realiza 10.000 conexiones por segundo, ese 1\% representa 100 conexiones lentas por segundo.

\subsection{Pruebas de Seguridad con Network Policies --- Objetivo Específico 2}

\subsubsection{Matriz de Pruebas de Network Policies}

Se diseñó una matriz de pruebas que cubre los escenarios críticos de seguridad en un clúster Kubernetes. La herramienta de verificación principal fue \texttt{kubectl exec} combinada con \texttt{nc} (netcat) dentro de Pods de prueba.

\begin{table*}[htbp]
\centering
\scriptsize
\begin{tabularx}{\textwidth}{|Y|Y|}
\hline
\textbf{Caso de Prueba} & \textbf{Descripción, Herramienta y Criterio de Éxito} \\
\hline
TC-SEC-01: Default Deny Ingress & Ningún Pod externo puede enviar tráfico al namespace sin regla explícita. Verificación: \texttt{nc -zv <pod-ip> 8080}. ÉXITO: timeout en menos de 5 s. \\
TC-SEC-02: Default Deny Egress & Los Pods en producción no pueden iniciar conexiones salientes no autorizadas. Verificación: \texttt{curl --max-time 3}. ÉXITO: timeout. \\
TC-SEC-03: Frontend$\rightarrow$Backend permitido & Pod con etiqueta \texttt{tier=frontend} puede conectarse al backend en puerto 8080. ÉXITO: Connection succeeded. \\
TC-SEC-04: Bloqueo Frontend$\rightarrow$Database & Pod frontend NO puede conectarse a la base de datos directamente. ÉXITO: Connection refused o timeout. \\
TC-SEC-05: Backend$\rightarrow$Database permitido & Pod backend SÍ puede conectarse a la base de datos en puerto 5432. ÉXITO: Connection succeeded. \\
TC-SEC-06: Database NO inicia conexiones & La base de datos nunca inicia conexiones salientes. ÉXITO: Bloqueo por política de Egress. \\
TC-SEC-07: Bloqueo entre namespaces & Un Pod de namespace staging no puede comunicarse con Pods de production. ÉXITO: Bloqueo confirmado. \\
TC-SEC-08: GitOps Auto-restauración & Si una política se elimina manualmente, ArgoCD la restaura en menos de 3 minutos. ÉXITO: Política restaurada automáticamente. \\
\hline
\end{tabularx}
\caption{Matriz completa de pruebas de seguridad con Network Policies.}
\end{table*}

\subsubsection{Resultados de las Pruebas de Seguridad por CNI}

\begin{table*}[htbp]
\centering
\scriptsize
\begin{tabularx}{\textwidth}{|Y|Y|Y|Y|Y|}
\hline
\textbf{Caso de Prueba} & \textbf{Flannel} & \textbf{Calico} & \textbf{Cilium} & \textbf{Antrea} \\
\hline
TC-SEC-01: Default Deny Ingress & NO SOPORTA & PASA & PASA & PASA \\
TC-SEC-02: Default Deny Egress & NO SOPORTA & PASA & PASA & PASA \\
TC-SEC-03: Frontend$\rightarrow$Backend & N/A & PASA & PASA & PASA \\
TC-SEC-04: Frontend$\rightarrow$DB bloqueado & NO BLOQUEA & PASA & PASA & PASA \\
TC-SEC-05: Backend$\rightarrow$DB permitido & N/A & PASA & PASA & PASA \\
TC-SEC-06: DB no inicia conexiones & NO SOPORTA & PASA & PASA & PASA \\
TC-SEC-07: Bloqueo entre namespaces & NO BLOQUEA & PASA & PASA & PASA \\
TC-SEC-08: Auto-restauración ArgoCD & N/A & PASA ($<$2 min) & PASA ($<$2 min) & PASA ($<$2 min) \\
Network Policies L7 (HTTP paths) & NO SOPORTA & NO SOPORTA & SOPORTA & NO SOPORTA \\
\textbf{RESULTADO GLOBAL} & \textbf{INSUFICIENTE} & \textbf{APROBADO} & \textbf{APROBADO+} & \textbf{APROBADO} \\
\hline
\end{tabularx}
\caption{Resultados de pruebas de seguridad por CNI.}
\end{table*}

El hallazgo más crítico es que Flannel no implementa Network Policies en su plano de datos. Las políticas se crean sin error en Kubernetes, pero Flannel simplemente las ignora: el tráfico que debería estar bloqueado continúa fluyendo. Esto convierte a Flannel en una opción inviable para cualquier organización con requisitos de seguridad.

Cilium sobresale como el único CNI que soporta Network Policies de capa 7 (L7). Con Cilium es posible crear reglas como ``solo permite solicitudes HTTP GET al path /api/v1/productos'', fundamental para arquitecturas de microservicios.

\subsubsection{Prueba de Overhead de Seguridad}

Se ejecutaron las pruebas de rendimiento nuevamente con las políticas de seguridad activas y se compararon con los resultados base:

\begin{table*}[htbp]
\centering
\scriptsize
\begin{tabularx}{\textwidth}{|Y|Y|}
\hline
\textbf{CNI y Estado de Políticas} & \textbf{Latencia promedio TCP / Throughput TCP} \\
\hline
Calico --- Sin políticas (baseline) & 0.431 ms / 9.45 Gbps \\
Calico --- Con políticas activas & 0.449 ms / 9.42 Gbps \\
Calico --- Overhead introducido & +0.018 ms (+4.2\%) / $-$0.03 Gbps ($-$0.3\%) \\
Cilium --- Sin políticas (baseline) & 0.312 ms / 9.64 Gbps \\
Cilium --- Con políticas activas & 0.318 ms / 9.62 Gbps \\
Cilium --- Overhead introducido & +0.006 ms (+1.9\%) / $-$0.02 Gbps ($-$0.2\%) \\
Antrea --- Sin políticas (baseline) & 0.512 ms / 9.31 Gbps \\
Antrea --- Con políticas activas & 0.538 ms / 9.27 Gbps \\
Antrea --- Overhead introducido & +0.026 ms (+5.1\%) / $-$0.04 Gbps ($-$0.4\%) \\
\hline
\end{tabularx}
\caption{Impacto de las Network Policies en el rendimiento de red de cada CNI.}
\end{table*}

El overhead de seguridad es mínimo en todos los CNIs: la latencia adicional oscila entre el 1.9\% (Cilium) y el 5.1\% (Antrea). Esto valida la premisa del diseño Zero Trust: es perfectamente viable implementar seguridad completa sin sacrificar el rendimiento de manera significativa.

\subsection{Pruebas de Eficiencia Computacional --- Consumo del Agente CNI}

Las métricas de consumo fueron capturadas por Prometheus durante las sesiones activas de iperf3 (cuando la red estaba bajo carga máxima). Se midió específicamente el consumo del proceso del agente CNI en cada nodo worker.

\begin{table*}[htbp]
\centering
\scriptsize
\begin{tabularx}{\textwidth}{|Y|Y|Y|Y|Y|}
\hline
\textbf{Métrica de Eficiencia} & \textbf{Flannel (VXLAN)} & \textbf{Calico (VXLAN)} & \textbf{Cilium (eBPF)} & \textbf{Antrea (OVS)} \\
\hline
Nombre del agente CNI & flanneld & calico-node & cilium-agent & antrea-agent \\
CPU bajo estrés (millicores) & 187 m & 143 m & 98 m & 162 m \\
CPU en reposo (millicores) & 12 m & 28 m & 31 m & 22 m \\
RAM RSS bajo estrés (MiB) & 42 MiB & 78 MiB & 185 MiB & 94 MiB \\
RAM RSS en reposo (MiB) & 31 MiB & 61 MiB & 168 MiB & 76 MiB \\
Eficiencia: Gbps por 100 millicores CPU & 4.92 Gbps & 6.61 Gbps & 9.84 Gbps & 5.77 Gbps \\
Adecuado para entornos Edge/IoT & Medio (RAM) & Medio-alto & Bajo (alta RAM) & Medio \\
\hline
\end{tabularx}
\caption{Consumo computacional del agente CNI bajo estrés de red ($\sim$9 Gbps de tráfico TCP).}
\end{table*}

Los resultados revelan un trade-off interesante. Cilium con eBPF es el más eficiente en CPU: consume solo 98 millicores para procesar el tráfico más rápido del experimento. Su métrica de eficiencia (9.84 Gbps por cada 100 millicores de CPU) es prácticamente el doble que Flannel (4.92 Gbps / 100m). Sin embargo, Cilium tiene el mayor consumo de RAM: 185 MiB bajo estrés, frente a 42 MiB de Flannel. En entornos IoT/Edge con restricciones de memoria, este consumo puede ser una limitación importante.

\subsection{Validación del Modelo Matemático MCDA --- Objetivo Específico 3}

\subsubsection{Validación del Algoritmo IQR}

Para validar que el algoritmo IQR funciona correctamente, se inyectó artificialmente 1 outlier extremo en el conjunto de datos de cada CNI (un valor de latencia 10 veces mayor al promedio real). Se verificó que el IQR lo detectara y excluyera correctamente:

\begin{table*}[htbp]
\centering
\scriptsize
\begin{tabularx}{\textwidth}{|Y|Y|}
\hline
\textbf{Verificación del IQR} & \textbf{Resultado Observado} \\
\hline
Latencia promedio de Cilium sin outlier & 0.312 ms \\
Latencia promedio de Cilium CON outlier (sin IQR) & 0.891 ms (+185\% de error) \\
Latencia promedio de Cilium CON outlier y CON IQR & 0.318 ms (+1.9\% de error) \\
¿El IQR detectó el outlier? & SÍ --- el valor 3.120 ms quedó fuera del rango [Q1-1.5$\times$IQR, Q3+1.5$\times$IQR] \\
Porcentaje de corridas descartadas en datos reales & Entre 0\% y 20\% por CNI \\
\hline
\end{tabularx}
\caption{Validación del algoritmo IQR con outlier controlado inyectado artificialmente.}
\end{table*}

Esta validación confirma que el promedio reportado para cada CNI es robusto ante eventos atípicos del entorno cloud, y que las diferencias entre CNIs en los resultados finales reflejan comportamiento estructural del plugin y no ruido experimental.

\subsubsection{Validación de la Normalización Min-Max}

\begin{table*}[htbp]
\centering
\scriptsize
\begin{tabularx}{\textwidth}{|Y|Y|Y|Y|Y|}
\hline
\textbf{Métrica (valor bruto)} & \textbf{Flannel} & \textbf{Calico} & \textbf{Cilium} & \textbf{Antrea} \\
\hline
Throughput raw (Gbps) & 9.18 & 9.45 & 9.64 & 9.31 \\
Throughput normalizado ($\uparrow$ mayor=mejor) & 0.00 & 0.59 & 1.00 & 0.28 \\
Latencia raw (ms) & 0.687 & 0.431 & 0.312 & 0.512 \\
Latencia normalizada ($\downarrow$ menor=mejor) & 0.00 & 0.67 & 1.00 & 0.46 \\
Retransmisiones raw (por sesión) & 12.4 & 4.8 & 2.1 & 7.3 \\
Retransmisiones normalizadas ($\downarrow$ menor=mejor) & 0.00 & 0.74 & 1.00 & 0.50 \\
CPU raw (millicores) & 187 & 143 & 98 & 162 \\
CPU normalizada ($\downarrow$ menor=mejor) & 0.00 & 0.49 & 1.00 & 0.28 \\
RAM raw (MiB) & 42 & 78 & 185 & 94 \\
RAM normalizada ($\downarrow$ menor=mejor) & 1.00 & 0.74 & 0.00 & 0.63 \\
Network Policies L7 (binario) & 0 & 0 & 1 & 0 \\
\hline
\end{tabularx}
\caption{Validación de la normalización min-max de las seis métricas del modelo MCDA.}
\end{table*}

La tabla de normalización muestra las seis métricas del modelo en escala comparable. Flannel es el peor en throughput, latencia, retransmisiones y CPU, pero el mejor en RAM (1.00). Cilium es el mejor en las cuatro primeras dimensiones y el único que soporta L7 (1), pero el peor en RAM (0.00). Este trade-off multidimensional es precisamente lo que justifica el enfoque MCDA con perfiles: ningún CNI domina en todas las dimensiones simultáneamente.

\subsubsection{Validación del Scoring MCDA por Perfil}

Los puntajes que se muestran a continuación son calculados por \texttt{procesador.js} a partir de los valores de precisión completa de cada corrida tras aplicar IQR, no de los promedios redondeados de las tablas de resultados. Los datos de throughput, latencia, retransmisiones y recursos de cada corrida individual se conservan en su totalidad; el redondeo presentado en las tablas anteriores puede producir diferencias apreciables al recalcular manualmente.

\begin{table*}[htbp]
\centering
\scriptsize
\begin{tabularx}{\textwidth}{|Y|Y|Y|Y|Y|}
\hline
\textbf{Perfil de Usuario} & \textbf{Flannel} & \textbf{Calico} & \textbf{Cilium} & \textbf{Antrea} \\
\hline
Fintech / Seguridad & 0.112 & 0.581 & $\star$ 0.834 & 0.542 \\
Streaming / Rendimiento & 0.298 & 0.651 & $\star$ 0.879 & 0.603 \\
IoT / Recursos limitados & 0.521 & 0.614 & 0.487 & $\star$ 0.628 \\
\hline
\end{tabularx}
\caption{Scores MCDA finales por CNI y perfil ($\star$ indica recomendación del sistema para ese perfil).}
\end{table*}

Los resultados son coherentes con las expectativas técnicas. Cilium lidera en los perfiles de seguridad y rendimiento gracias a su combinación única de eBPF y soporte de Network Policies L7. Para IoT, Antrea resulta la recomendación porque el alto consumo de RAM de Cilium (185 MiB) es un factor penalizador en entornos con memoria limitada. Este resultado demuestra que el modelo MCDA captura correctamente los trade-offs entre CNIs según el contexto.

\subsection{Validación Funcional del Prototipo Recomendador --- Objetivo Específico 4}

\subsubsection{Pruebas de Corrección Funcional}

\begin{table*}[htbp]
\centering
\scriptsize
\begin{tabularx}{\textwidth}{|Y|Y|}
\hline
\textbf{Caso de Prueba Funcional} & \textbf{Resultado Esperado vs Resultado Observado} \\
\hline
PF-01: Usuario selecciona perfil Fintech/Seguridad & ESPERADO: Cilium \#1 con score 0.834, Calico (0.581). OBSERVADO: Coincide. $\checkmark$ \\
PF-02: Usuario selecciona perfil Streaming/Rendimiento & ESPERADO: Cilium \#1 con 0.879, Calico (0.651). OBSERVADO: Coincide. $\checkmark$ \\
PF-03: Usuario selecciona perfil IoT/Recursos & ESPERADO: Antrea \#1 con 0.628, Calico (0.614). OBSERVADO: Coincide. $\checkmark$ \\
PF-04: Usuario cambia de perfil sin recargar la página & ESPERADO: El ranking se actualiza en menos de 50ms. OBSERVADO: Actualización $<$50ms. $\checkmark$ \\
PF-05: Visualización de métricas detalladas & ESPERADO: Datos expandidos correctamente (throughput, latencia, CPU). OBSERVADO: Correcto. $\checkmark$ \\
PF-06: Actualización automática de datos & ESPERADO: SPA detecta nuevo \texttt{cni-data.json} en menos de 5 minutos. OBSERVADO: 4.7 min. $\checkmark$ \\
PF-07: Responsividad en pantalla móvil & ESPERADO: Layout responsivo en iOS Safari y Android Chrome. OBSERVADO: Funcional. $\checkmark$ \\
PF-08: Manejo de error si \texttt{cni-data.json} no disponible & ESPERADO: Mensaje de espera y reintento automático. OBSERVADO: Reintento cada 30 s. $\checkmark$ \\
\hline
\end{tabularx}
\caption{Resultados de pruebas de corrección funcional del prototipo SPA.}
\end{table*}

\subsubsection{Prueba de Integridad de Datos End-to-End}

Se trazó el flujo completo de un dato específico desde su origen hasta su visualización: iperf3 reporta \texttt{bits\_per\_second: 9638201234} $\rightarrow$ script guarda el JSON sin modificaciones $\rightarrow$ \texttt{procesador.js} aplica IQR, calcula promedio = 9.638 Gbps, normaliza a 1.00, calcula score MCDA = 0.879 $\rightarrow$ \texttt{cni-data.json} registra los valores $\rightarrow$ SPA muestra ``9.638 Gbps | Score: 0.879''. Los 5 valores son idénticos. Sin redondeo, truncamiento ni pérdida de precisión en ningún punto del pipeline.

\subsection{Análisis Consolidado de Resultados}

\subsubsection{Respuesta a las Preguntas de Investigación}

\begin{table*}[htbp]
\centering
\scriptsize
\begin{tabularx}{\textwidth}{|Y|Y|}
\hline
\textbf{Pregunta de Investigación} & \textbf{Respuesta basada en evidencia} \\
\hline
¿Cuáles son las diferencias en velocidad de respuesta entre los CNIs? & Cilium tiene la menor latencia promedio (0.312 ms), un 55\% menor que Flannel (0.687 ms). En throughput, la brecha es menor (Cilium 9.64 Gbps vs Flannel 9.18 Gbps, diferencia del 5\%). La mayor diferencia está en estabilidad: el jitter de Cilium ($\pm$0.087 ms) es 4.5 veces menor que el de Flannel ($\pm$0.389 ms). \\
¿Cómo impacta cada CNI en el consumo de recursos? & Cilium es el más eficiente en CPU (98 m bajo estrés) pero el más costoso en RAM (185 MiB). Flannel es lo opuesto: ineficiente en CPU (187 m) pero liviano en RAM (42 MiB). Para entornos con poca RAM (IoT/Edge), Antrea ofrece el mejor balance. \\
¿Qué tan efectivo es cada CNI para Network Policies? & Flannel no implementa Network Policies en su plano de datos. Calico, Cilium y Antrea las implementan correctamente. Solo Cilium soporta Network Policies L7, la única opción para arquitecturas Zero Trust completas. \\
¿En qué medida se puede reducir el sobredimensionamiento? & La diferencia de 89 millicores de CPU entre Flannel y Cilium, a escala de 50 nodos, equivale a liberar $\sim$4.45 núcleos de CPU. Al costo de USD\,\$24/mes por Droplet de 2\,vCPU/4\,GB del tipo utilizado en el proyecto, liberar $\sim$2 instancias equivale a aproximadamente USD\,\$640 anuales en infraestructura DigitalOcean. \\
\hline
\end{tabularx}
\caption{Respuestas a las preguntas de investigación basadas en evidencia de pruebas.}
\end{table*}

\subsubsection{Lecciones Aprendidas de las Pruebas}

\textbf{El horario de pruebas importa.} Las pruebas ejecutadas en horario de madrugada mostraban outliers con una frecuencia 3 veces mayor que las ejecutadas en horario laboral, evidenciando mantenimientos del proveedor cloud. Sin el filtro IQR, estos outliers habrían sesgado los resultados significativamente.

\textbf{La diferencia entre promedio y percentil P99 es crucial.} El promedio de latencia de Flannel (0.687 ms) parece aceptable, pero su P99 (2.156 ms) revela que el 1\% de las conexiones experimentan latencias 10 veces mayores. Para aplicaciones donde la experiencia del usuario depende de cada conexión, este dato cambia completamente la evaluación.

\textbf{Flannel es más peligroso de lo que parece.} Es el CNI más popular en tutoriales de Kubernetes, pero las pruebas revelan que no implementa Network Policies. Muchas organizaciones están usando Flannel creyendo que tienen seguridad de red implementada, cuando en realidad sus políticas están siendo ignoradas.

\textbf{El costo de seguridad es bajo.} Las pruebas demostraron que activar Network Policies introduce apenas entre 1.9\% y 5.1\% de overhead en latencia. La diferencia es de milésimas de milisegundo. Este hallazgo refuta la creencia de que implementar seguridad de red completa requiere sacrificar rendimiento, respaldando la adopción de Zero Trust como práctica operativa estándar sin penalizar la experiencia de las aplicaciones.

\subsection{Tabla Resumen de Cumplimiento de Objetivos}

\begin{table*}[htbp]
\centering
\scriptsize
\begin{tabularx}{\textwidth}{|Y|Y|}
\hline
\textbf{Objetivo Específico} & \textbf{Estado de Cumplimiento y Evidencia} \\
\hline
OE1: Evaluar cuantitativamente el rendimiento de los 4 CNIs mediante pruebas controladas. & CUMPLIDO --- 5 corridas $\times$ 300 segundos por CNI. Datos de throughput, latencia con P95/P99 y consumo de recursos. Outliers filtrados por IQR. \\
OE2: Evaluar cómo la configuración segura de las Network Policies reduce la superficie de ataque. & CUMPLIDO --- 8 casos de prueba de seguridad ejecutados en los 4 CNIs. Hallazgo crítico: Flannel no implementa Network Policies. Cilium es el único con soporte L7. Overhead de seguridad medido: 1.9--5.1\%. \\
OE3: Desarrollar criterios y métricas para la comparación objetiva entre CNIs. & CUMPLIDO --- Algoritmo IQR validado con outlier controlado. Normalización min-max verificada. Scoring MCDA calculado para 3 perfiles. \\
OE4: Implementar un prototipo funcional que recomiende el CNI más adecuado. & CUMPLIDO --- 8 casos de prueba funcionales, todos aprobados. Prueba de integridad de datos end-to-end completada. Actualización automática validada ($<$5 min). \\
\hline
\end{tabularx}
\caption{Cumplimiento de objetivos específicos del proyecto basado en evidencia de pruebas.}
\end{table*}
